\chapter{最小 CPU 与控制 trace}

前面已经具备构造处理器的必要部件：组合逻辑、状态元件、时钟边界、ALU、数据通路和控制器。本章只做一件事：把这些部件合成为一台能够取指、译码、执行、写回和更新 PC 的最小 CPU。主存、总线、设备和经典芯片会在后续章节展开，本章先把 CPU 内部的状态边界和控制 trace 写清楚。

CPU 可以视为一组同步硬件闭环：PC 指出下一条编码地址，取指路径取回指令 bit，译码逻辑生成控制信号，数据通路读取和计算，状态元件在时钟边界提交结果。只要能为一条指令写出 PC、IR、寄存器堆、ALU、写回 mux、标志位和控制器的完整 trace，就已经能判断这台最小 CPU 是否真的会执行程序。

\begin{keyidea}
最小 CPU 不是“ALU 加寄存器”的粗线图，而是一组按周期提交的状态转换。每条指令都必须说明：本周期读哪些旧状态，组合路径经过哪些部件，哪些控制线有效，时钟沿会提交哪些新状态，异常或等待是否会阻止提交。
\end{keyidea}

\begin{pdfdiagram}
  \node[hw state, minimum width=1.7cm] (pc) at (0,0.75) {PC};
  \node[hw compute, minimum width=2.0cm] (imem) at (2.55,0.75) {指令存储};
  \node[hw state, minimum width=1.55cm] (ir) at (5.05,0.75) {IR};
  \node[hw control, minimum width=1.9cm] (ctrl) at (7.55,0.75) {控制器};
  \node[hw state, minimum width=2.0cm] (rf) at (2.55,-1.0) {寄存器堆};
  \node[hw compute, minimum width=1.7cm] (alu) at (5.05,-1.0) {ALU};
  \node[hw compute, minimum width=2.0cm] (wb) at (7.55,-1.0) {写回 mux};
  \draw[hw bus] (pc) -- node[above=2pt,hw label,fill=white,inner sep=1pt]{取指地址} (imem);
  \draw[hw bus] (imem) -- node[above=2pt,hw label,fill=white,inner sep=1pt]{编码 bit} (ir);
  \draw[hw bus] (ir) -- node[above=2pt,hw label,fill=white,inner sep=1pt]{字段} (ctrl);
  \draw[hw return] (ctrl.south) |- node[pos=0.25,right,hw label,fill=white,inner sep=1pt]{控制线} (5.05,-0.18) -| (alu.north);
  \draw[hw bus] (rf) -- node[above=2pt,hw label,fill=white,inner sep=1pt]{源操作数} (alu);
  \draw[hw bus] (alu) -- node[above=2pt,hw label,fill=white,inner sep=1pt]{结果/标志} (wb);
  \draw[hw return] (wb.east) -- ++(0.55,0) |- (rf.south) node[pos=0.25,below=2pt,hw label,fill=white,inner sep=1pt]{写回};
  \draw[hw return] (ctrl.north) |- (3.85,1.75) -| (pc.north) node[pos=0.2,above=2pt,hw label,fill=white,inner sep=1pt]{下一 PC};
  \node[hw label, text width=9.4cm] at (4.1,-2.05) {最小 CPU 的核心验收对象是 trace：PC 取指，IR 保存编码，控制器打开数据通路，寄存器和 PC 在时钟沿提交。};
\end{pdfdiagram}
\diagramnote{最小 CPU 总图。箭头表示一次指令从取指到写回的证据链，反馈线表示 PC 和寄存器堆在时钟边界更新。}

\section{一、从动作编码到最小 CPU}

若希望同一套数据通路在不同时间做不同动作，就需要把“这一次做什么”也表示成 bit。动作编码可以先写成最小形式：

\begin{lstlisting}
opcode | dst | src_a | src_b
\end{lstlisting}

\code{opcode} 表示动作类型，例如 ADD、SUB、AND；\code{dst} 表示写入哪个寄存器；\code{src\_a} 和 \code{src\_b} 表示读取哪两个寄存器。控制器读取这些字段后，把它们译成 ALU op、读地址、写地址、写使能和写回选择。编码本身仍然是硬件输入，不是文字命令。

若有一组动作编码连续放在存储结构中，硬件还需要知道当前取哪一条。这个位置由 PC 保存。PC 本质上也是寄存器，只是它保存的是“下一条动作编码的地址”。

\begin{longtable}{L{0.24\linewidth}L{0.52\linewidth}}
\toprule
部件 & 作用 \\
\midrule
PC & 保存下一条编码地址 \\
指令存储器 & 根据 PC 输出当前编码 \\
控制器 & 根据编码字段产生控制信号 \\
寄存器阵列 & 保存通用数据 \\
ALU & 做整数和逻辑运算 \\
选择器网络 & 选择 ALU 输入、写回来源和下一 PC \\
\bottomrule
\end{longtable}

顺序执行时，下一 PC 可以由加法器产生：

\begin{lstlisting}
pc_next = pc_current + instruction_size
\end{lstlisting}

如果每条编码占 4 个字节，\code{instruction\_size} 就是 4；如果是可变长度指令，下一 PC 的计算会更复杂。关键事实是：顺序执行不是抽象的“下一行”，而是 PC 这个状态寄存器保存了新的地址。

最小单周期模型可以这样走：

\begin{longtable}{L{0.18\linewidth}L{0.52\linewidth}L{0.18\linewidth}}
\toprule
阶段 & 硬件动作 & 边界 \\
\midrule
取指 & PC 送入指令存储器，得到编码 & 组合路径 \\
译码 & 控制器产生 mux、ALU、写使能信号 & 组合路径 \\
执行 & 寄存器阵列读值，ALU 计算 & 组合路径 \\
写回 & 结果在时钟沿写入寄存器或 PC & 状态边界 \\
\bottomrule
\end{longtable}

用一条普通运算编码走读：当前位置 PC=100，指令存储器在地址 100 处保存的是：

\begin{lstlisting}
ADD dst=R3, src_a=R1, src_b=R2
\end{lstlisting}

周期开始时，PC 输出稳定为 100。这个地址进入指令存储器，存储器经过地址译码后把这条编码输出。控制器看到 opcode=ADD，于是把 \code{alu\_op} 设为 ADD，把寄存器阵列的两个读地址分别设为 R1 和 R2，把写地址设为 R3，把 \code{wb\_sel} 设为 ALU，把 \code{reg\_write} 设为 1，同时让下一 PC 选择 PC 加指令长度。

这些控制信号稳定后，寄存器阵列的两个读端口输出 R1 和 R2 当前的值。它们经过 ALU 输入选择器进入 ALU，加法路径产生结果。结果再经过写回选择器，到达寄存器阵列的写数据端口。另一边，PC 加法器准备下一条指令地址，PC 选择器把它送到 PC 的下一值输入。只有到时钟沿，两个状态变化才同时提交：R3 保存 ALU 结果，PC 保存下一条地址。

分支也可以归结为下一 PC 选择。顺序情况下，下一 PC 是 PC 加指令长度；条件分支则由 ALU 比较结果和选择器决定下一 PC 来源：

\begin{lstlisting}
branch_taken = condition_met AND branch_op
next_pc = mux(branch_taken, pc_sequential, branch_target)
\end{lstlisting}

分支的硬件本质是改变 PC 这个寄存器的下一值。下一周期取到哪条编码，完全由这个新 PC 决定。

\topic{把指令周期写成可验收账本}经典教材讲 CPU 时，常把一条指令拆成取指、译码、执行、访存、写回和下一 PC 更新。这个拆分不是为了背阶段名称，而是为了把每个阶段对应到真实硬件路径。最小单周期 CPU 可以在一个长周期内完成所有路径；多周期 CPU 可以把这些路径分到多个时钟；流水线 CPU 可以让多条指令在不同阶段重叠推进。三者的共同点是：每个阶段都要说明源状态、组合路径、目标状态和异常边界。

\begin{longtable}{L{0.14\linewidth}L{0.34\linewidth}L{0.28\linewidth}L{0.14\linewidth}}
\toprule
阶段 & 主要硬件动作 & 关键控制信号 & 提交边界 \\
\midrule
IF 取指 & PC 送到指令存储器或 I-cache，取回指令 bit & \code{pc_out}、\code{imem_read}、\code{ir_write} & IR 或取指缓冲 \\
ID 译码 & opcode、寄存器号、立即数字段被解释成控制线 & \code{decode}、\code{imm_kind}、\code{read_addr} & 可组合，也可锁存 \\
EX 执行 & ALU 做加减、逻辑、比较或地址形成 & \code{alu_op}、\code{a_sel}、\code{b_sel} & 输出和标志 \\
MEM 访存 & 对数据存储器、cache 或 MMIO 发起读写事务 & \code{mem_read}、\code{mem_write}、\code{byte_en} & MDR 或外部响应 \\
WB 写回 & ALU 结果或读回数据写入寄存器堆 & \code{wb_sel}、\code{reg_write} & 目标寄存器 \\
PC 更新 & 选择顺序 PC、分支目标、跳转目标或异常入口 & \code{pc_sel}、\code{pc_write} & PC \\
\bottomrule
\end{longtable}

这张表可以直接变成验收清单。若某条 LOAD 指令读到了错误数据，不应只问“内存为什么错”，而要逐项检查：PC 是否取到正确编码，译码是否识别为 LOAD，ALU 是否形成正确地址，MEM 阶段是否访问正确目标，读回数据是否进入 MDR，WB 阶段是否选择了 MDR 而不是 ALUOut。这样的排查方式比凭感觉改控制器可靠得多。

\topic{控制字把“会执行指令”落到可接线信号}教学 CPU 最容易写成口号：取指、译码、执行、写回。真正能搭出来的版本必须有控制字。控制字不是高级软件结构，而是一组在某个周期稳定的控制线集合。它决定寄存器是否写入、ALU 选哪个操作、存储器是否读写、PC 从哪里取下一值。若控制字没有定义清楚，CPU 即使图上连通，也无法稳定执行程序。

\begin{longtable}{L{0.18\linewidth}L{0.2\linewidth}L{0.2\linewidth}L{0.2\linewidth}L{0.12\linewidth}}
\toprule
控制对象 & ADD & LOAD & STORE & JZ \\
\midrule
ALU 操作 & 加法或指定算术 & 基址加偏移 & 基址加偏移 & 比较或测试零 \\
寄存器读 & 两个源寄存器 & 基址寄存器 & 基址和数据寄存器 & 条件源寄存器 \\
存储器读 & 无 & 有 & 无 & 无 \\
存储器写 & 无 & 无 & 有 & 无 \\
写回选择 & ALU 结果 & 存储器数据 & 无 & 无 \\
PC 选择 & 顺序 PC & 顺序 PC & 顺序 PC & 条件目标或顺序 PC \\
\bottomrule
\end{longtable}

若把这些控制对象压成 bit，可以得到类似下面的教学控制字。具体 bit 位不重要，重要的是每一位都能对应到实际连线或选择器输入。

\begin{lstlisting}
control word fields:
  pc_sel      // sequential, branch_target, jump_target, exception
  reg_write   // enable register-file write port
  wb_sel      // ALU, memory, pc_plus_size
  alu_op      // add, sub, and, or, compare
  alu_b_sel   // register or immediate
  mem_read    // request data read
  mem_write   // request data write
  byte_en     // selected byte lanes
  flags_write // update Zero/Carry/Overflow/Sign
\end{lstlisting}

\topic{硬连线控制与微程序控制是两种组织控制字的方法}控制器可以直接由 opcode、状态位和当前周期译码出控制线，这称为硬连线控制。它速度快，结构适合简单指令集和高频流水线，但修改指令行为需要改组合逻辑。另一种方法是把控制字存在控制存储器中，由微程序计数器逐条取出微指令；每条机器指令对应一段微操作序列。这种方法更像“硬件内部的小程序”，适合组织复杂指令、多周期总线访问和异常入口，但控制存储器和微序列会增加延迟。

\begin{longtable}{L{0.18\linewidth}L{0.34\linewidth}L{0.38\linewidth}}
\toprule
控制方式 & 适合场景 & 验收问题 \\
\midrule
硬连线控制 & 指令格式规整、周期少、目标频率高 & opcode、状态位和时序状态是否唯一决定控制线 \\
微程序控制 & 指令复杂、多周期动作多、需要复用数据通路 & 微地址如何选择、分支微指令如何跳转、异常如何打断 \\
混合方式 & 常见简单指令走快路径，复杂指令走微序列 & 快路径和微序列是否共享一致的提交边界 \\
\bottomrule
\end{longtable}

\begin{lstlisting}
microinstruction sketch:
  alu_op | src_a | src_b | dst | mem_cmd | pc_cmd | next_micro_pc

LOAD sequence:
  u0: IR  <- MEM[PC],       next = u1
  u1: A   <- R[base],       next = u2
  u2: ADR <- A + imm,       next = u3
  u3: MDR <- MEM[ADR],      next = u4 or wait
  u4: R[dst] <- MDR, PC <- PC+4, next = fetch
\end{lstlisting}

微程序写法有一个重要好处：它迫使作者说明每个周期保存什么中间状态。若某一步需要等待外部存储器，\code{next\_micro\_pc} 可以停留在当前微地址；若访问失败，微序列可以跳到异常入口。这样，等待、错误和普通写回不再混在一句“执行 LOAD”里。

\topic{多周期 CPU 用微操作缩短最长组合路径}单周期 CPU 概念清楚，但周期必须长到容纳取指、译码、ALU、访存和写回的最坏组合路径。多周期 CPU 把同一条指令拆成若干微操作，让每个周期只完成一小段可验收动作。例如 LOAD 可以拆为：PC 取指、译码读寄存器、ALU 形成地址、存储器读、写回寄存器。每一步之间用 IR、MDR、ALUOut 或临时寄存器保存中间结果。

\begin{longtable}{L{0.12\linewidth}L{0.28\linewidth}L{0.3\linewidth}L{0.2\linewidth}}
\toprule
周期 & 微操作 & 打开的硬件路径 & 提交对象 \\
\midrule
T0 & \code{IR = MEM[PC]} & PC 到指令存储器，读响应到 IR & IR \\
T1 & \code{A = R[rs], B = R[rt]} & 指令字段到寄存器堆读地址 & A/B 临时寄存器 \\
T2 & \code{ALUOut = A + imm} & ALU 做地址形成或比较 & ALUOut/Flags \\
T3 & \code{MDR = MEM[ALUOut]} & 地址到数据存储器，读响应到 MDR & MDR \\
T4 & \code{R[rd] = MDR, PC = PC+4} & 写回 mux 和 PC 加法器 & 寄存器、PC \\
\bottomrule
\end{longtable}

这个 trace 说明，多周期设计不是“慢版单周期”，而是把过长的硬件路径切成可计时、可观察、可插入等待的边界。若数据存储器慢，可以只让 T3 等待；若分支比较已经在 T2 得到结果，可以在 T2 或 T3 改写 PC；若发生异常，可以把 PC 选择器改到异常入口，而不是继续提交错误写回。

\topic{先规定一个极小 ISA，程序才能进入 ROM}为了让整机实验真正运行，需要把“指令”压成固定 bit。下面是一种 16-bit 教学编码：高 4 bit 是 opcode，中间字段选择寄存器，低位可解释为寄存器号、立即数或分支偏移。它不是任何真实 ISA 的复制品，但足以把编码、译码、控制字和 ROM 内容连接起来。

\begin{longtable}{L{0.16\linewidth}L{0.22\linewidth}L{0.24\linewidth}L{0.28\linewidth}}
\toprule
指令 & 编码格式 & 语义 & 需要的硬件路径 \\
\midrule
\code{ADD} & \code{0001 d a b ----} & \code{R[d]=R[a]+R[b]} & 双读端口、ALU、写回 \\
\code{LDI} & \code{0010 d imm8} & \code{R[d]=zero_extend(imm8)} & 立即数扩展、写回 \\
\code{LOAD} & \code{0011 d a off4} & \code{R[d]=MEM[R[a]+off]} & 地址形成、读存储、写回 \\
\code{STORE} & \code{0100 s a off4} & \code{MEM[R[a]+off]=R[s]} & 地址形成、写数据、字节使能 \\
\code{JZ} & \code{0101 r off8} & 若 \code{R[r]==0} 则 PC 加偏移 & 比较、符号扩展、PC 选择 \\
\code{JMP} & \code{0110 off12} & PC 加偏移 & 目标形成、PC 写入 \\
\bottomrule
\end{longtable}

这个编码能直接写入 ROM。若 \code{GPIO_IN=0xF004}、\code{GPIO_OUT=0xF000}，程序可以先用 \code{LDI} 装入高地址基址，再循环读取按键、条件跳转、写 LED。教学实现中可把 \code{R0} 固定为 0，也可用一个寄存器保存 I/O 基址。关键不是汇编语法，而是 ROM 中每个 16-bit 字都能被译码成控制信号。

\begin{lstlisting}
tiny ROM program sketch:
  LDI   R1, 0xF0       // R1 holds high part of MMIO base
loop:
  LOAD  R2, [R1 + 4]   // read GPIO_IN
  JZ    R2, off
  LDI   R3, 1
  STORE R3, [R1 + 0]   // write GPIO_OUT = 1
  JMP   loop
off:
  LDI   R3, 0
  STORE R3, [R1 + 0]   // write GPIO_OUT = 0
  JMP   loop
\end{lstlisting}

\topic{程序员可见状态不是全部硬件状态}
机器级程序通常只直接看见一小组状态：PC、通用寄存器、条件码或标志、可寻址内存，有些 ISA 还显式暴露栈指针、链接寄存器或状态寄存器。CPU 内部还会有 IR、MDR、ALUOut、流水线寄存器、旁路网络和微程序计数器，但这些通常不是普通程序可以直接读写的对象。区分这两类状态很重要：ISA 规定的是程序必须可依赖的边界，微结构内部状态则服务于实现速度、面积和功耗。

\begin{longtable}{L{0.18\linewidth}L{0.36\linewidth}L{0.36\linewidth}}
\toprule
状态 & 程序员视角 & 硬件实现视角 \\
\midrule
PC & 下一条或当前执行位置 & 取指地址、分支选择、异常入口和返回路径 \\
通用寄存器 & 保存整数、地址、临时值和返回值 & 寄存器堆端口、旁路、写回优先级 \\
条件码 & 最近运算留下的比较证据 & ALU 标志生成、\code{flags_write} 和分支条件译码 \\
内存 & 按地址访问的字节对象 & cache、TLB、总线、MMIO 和权限检查 \\
栈指针 & 当前栈顶位置 & 普通寄存器加受约束的加减和 LOAD/STORE 序列 \\
\bottomrule
\end{longtable}

这种划分能解释为什么同一套 ISA 可以有不同实现。一个简单教学 CPU 可以用多周期状态机执行 LOAD；一个高性能 CPU 可以把它拆进流水线、cache 和乱序队列。只要最终提交到程序员可见状态的结果相同，内部路径可以不同。反过来，若内部优化改变了 PC、寄存器、内存或条件码的可见顺序，就不再只是实现细节，而是破坏了 ISA 合同。

\topic{条件码把 ALU 证据变成控制流}
条件分支不是“看高级语言的 if”，而是看硬件已经保存的证据。ALU 做加法、减法或比较时，可以同时产生 Zero、Sign、Carry、Overflow 等标志。控制器把这些标志与指令中的条件字段组合，决定下一 PC 选择顺序地址还是分支目标。比较指令常常等价于一次不写回结果的减法：结果本身丢弃，但标志被保留下来供后续分支使用。

\begin{longtable}{L{0.16\linewidth}L{0.34\linewidth}L{0.4\linewidth}}
\toprule
条件 & 常见硬件证据 & 容易犯错的解释 \\
\midrule
相等 & \code{A-B} 的 Zero 为 1 & 只适合判断 bit 模式相等，不说明大小 \\
无符号小于 & 减法借位或 Carry 约定 & 不能直接用结果最高位判断 \\
有符号小于 & Sign 与 Overflow 的组合 & 最高位受溢出影响，必须修正 \\
结果为负 & Sign 标志为 1 & 只在补码解释下表示负数 \\
有符号溢出 & 同号相加得到异号，或减法符号规则触发 & 与无符号进位不是同一件事 \\
\bottomrule
\end{longtable}

\begin{lstlisting}
CMP R1, R2       // flags = R1 - R2, result is not written
JE  equal_path   // branch if Zero == 1
JL  less_path    // signed branch uses Sign xor Overflow
JB  below_path   // unsigned branch uses borrow/carry convention
\end{lstlisting}

这个例子说明，分支条件必须绑定解释规则。同一对寄存器 \code{R1} 和 \code{R2} 的 bit 模式，按无符号数、有符号补码数或地址偏移解释，可能得到不同的“小于”。硬件不会自动知道程序员想要哪一种比较，必须由 opcode 或条件字段告诉控制器使用哪组标志逻辑。

\topic{调用、返回和栈帧是受约束的地址事务}
函数调用可以从软件语法降到三件硬件事：保存返回地址，跳到被调用函数入口，为局部对象和保存寄存器预留内存。返回时再恢复必要状态，把 PC 改回返回地址。很多 ISA 用专门的 CALL/RET 指令完成其中一部分；精简教学 CPU 也可以用普通 STORE、LOAD、加减栈指针和 JUMP 组合出来。

\begin{lstlisting}
CALL target, stack grows downward:
  SP = SP - word_size
  MEM[SP] = PC + instruction_size   // push return address
  PC = target

RET:
  PC = MEM[SP]                       // pop return address
  SP = SP + word_size
\end{lstlisting}

\begin{longtable}{L{0.18\linewidth}L{0.38\linewidth}L{0.34\linewidth}}
\toprule
栈帧对象 & 机器级表示 & 硬件动作 \\
\midrule
返回地址 & 一个位宽等于 PC 的地址值 & STORE 到栈，RET 时 LOAD 到 PC \\
保存寄存器 & 调用约定规定要保留的寄存器内容 & 进入时 STORE，返回前 LOAD \\
局部变量 & SP 或帧指针加固定偏移 & 地址形成后普通 LOAD/STORE \\
实参溢出区 & 寄存器放不下的参数位置 & 调用方或被调用方按约定访问 \\
对齐空洞 & 为总线宽度、ABI 或向量访问保留的字节 & 调整 SP，避免未对齐访问代价或异常 \\
\bottomrule
\end{longtable}

把调用栈看成地址事务后，许多系统错误会变得具体。栈越界写不是抽象的“内存不安全”，而是一次 STORE 形成了错误地址，可能覆盖相邻局部变量、保存寄存器、返回地址或未映射页。若返回地址被覆盖，RET 只是按合同把栈顶字节装进 PC；CPU 并不知道这个地址原本是否可信。后续章节讨论保护模式、页表、异常和执行权限时，就是在给这些地址事务加上硬件可检查的边界。

\topic{从一条语句贯通到硬件事务}为了避免本章变成若干概念并列，可以用一条简单语句贯通整机：\code{if (MMIO[GPIO_IN] != 0) MMIO[GPIO_OUT] = 1;}。在软件层面它像一次条件判断；在硬件层面，它至少包含取指、译码、地址形成、MMIO 读、条件分支、MMIO 写和 PC 更新。每一步都可以写成可验收事务。

\begin{longtable}{L{0.16\linewidth}L{0.38\linewidth}L{0.36\linewidth}}
\toprule
层次 & 发生的动作 & 必须检查的证据 \\
\midrule
取指 & PC 指向 ROM 或 I-cache 中的 LOAD/JZ/STORE 编码 & PC、指令存储片选、IR 内容 \\
译码 & opcode 被译成 \code{mem_read}、\code{pc_sel}、\code{mem_write} 等控制线 & 控制字和寄存器读地址 \\
地址形成 & 基址寄存器与偏移形成 \code{GPIO_IN} 或 \code{GPIO_OUT} 地址 & ALU 输入、ALUOut、地址总线 \\
MMIO 读 & GPIO 控制器返回已同步按键状态 & GPIO 片选、读使能、返回数据稳定 \\
条件判断 & 比较结果决定顺序 PC 或分支目标 & Zero 标志、分支控制、下一 PC \\
MMIO 写 & 输出锁存器接收写数据并驱动 LED & 写使能、字节使能、锁存器输出 \\
\bottomrule
\end{longtable}

这个例子也说明，编译后的指令顺序与外设副作用不能随意重排。普通内存访问可以被 cache、写缓冲和预取优化；但 \code{GPIO_OUT} 写入是外部可见动作，必须按设备属性和顺序规则执行。若系统把 MMIO 当作普通 cacheable 内存，LED 可能不会在预期时刻改变，甚至写入被合并或延迟到错误位置。

\topic{五类基本指令覆盖了 CPU 的主要通路}最小 CPU 不必一开始支持复杂指令，但至少要能解释五类动作：寄存器到寄存器运算、加载、存储、条件分支和无条件跳转。它们共同覆盖寄存器堆、ALU、数据存储器、PC 选择器和标志位。

\begin{longtable}{L{0.14\linewidth}L{0.34\linewidth}L{0.32\linewidth}L{0.1\linewidth}}
\toprule
指令类 & 典型路径 & 必须打开的控制 & 写回对象 \\
\midrule
ADD/SUB & 寄存器读出，经 ALU，结果写回寄存器 & \code{alu_op}、\code{wb_sel=ALU}、\code{reg_write} & 寄存器、Flags \\
LOAD & 基址寄存器加立即数形成地址，存储器读回数据 & \code{alu_op=ADD}、\code{mem_read}、\code{wb_sel=MEM} & 寄存器 \\
STORE & 基址寄存器加立即数形成地址，源寄存器写入存储器 & \code{alu_op=ADD}、\code{mem_write}、\code{store_data_sel} & 存储器或 MMIO \\
BRANCH & 比较源操作数或标志位，条件成立则选择分支目标 & \code{alu_op=CMP}、\code{pc_sel}、\code{pc_write} & PC \\
JUMP & 立即数、寄存器或链接地址决定下一 PC & \code{target_sel}、\code{pc_sel}、可选 \code{reg_write} & PC，可选链接寄存器 \\
\bottomrule
\end{longtable}

以 LOAD 为例，\code{R3 = MEM[R1 + 8]} 不是单纯“读内存”。硬件要先读出 R1，把立即数 8 符号扩展或零扩展到 ALU 输入宽度，ALU 做地址加法，地址进入数据存储路径，读响应回来后经过写回 mux，最后在时钟沿写入 R3。若目标地址落入普通 RAM，读回的是存储阵列内容；若目标地址落入 MMIO 区，读回的是设备寄存器当前状态。地址相同形式，语义由地址译码决定。

\begin{lstlisting}
LOAD R3, [R1 + 8]
  EX:  addr = R1 + sign_extend(8)
  MEM: read target selected by addr
  WB:  R3 = read_data
\end{lstlisting}

STORE 的风险更高，因为它有副作用。 \code{MEM[R1 + 8] = R3} 若落到 RAM，表示保存数据；若落到设备命令寄存器，可能启动设备、清除中断或改变输出引脚。因此 STORE 必须明确字节使能、写数据来源、地址空间属性和完成条件。普通内存写响应表示写事务被接受；设备命令写响应通常不表示设备动作完成。

\begin{lstlisting}
STORE [R1 + 8], R3
  EX:  addr = R1 + sign_extend(8)
  MEM: write_data = R3, byte_en = access_size
       target receives write command
  WB:  no register writeback
\end{lstlisting}

BRANCH 和 JUMP 则说明 PC 是数据通路的一部分。条件分支一般要先得到比较证据，例如 Zero、Sign、Carry 或 Overflow，再由控制器决定 \code{pc_sel}。无条件跳转可以直接选择目标地址；带链接跳转还会把返回地址写入寄存器。此时同一条指令可能同时写 PC 和一个通用寄存器，控制器必须明确两个写边界是否在同一时钟沿提交。

\topic{基本块和循环把 PC 更新变成可读结构}
机器级程序可以先按基本块理解。一个基本块内部只有顺序执行，入口只有一个，出口通常是分支、跳转、返回或异常边界。硬件不认识“基本块”这个软件概念，但基本块对读 trace 很有用：块内每条指令都让 PC 顺序前进，块尾那条指令才选择下一 PC。循环就是某个块尾分支把 PC 改回较低地址。

\begin{lstlisting}
loop:
  LOAD R2, [R1 + 0]
  ADD  R3, R3, R2
  ADDI R1, R1, 4
  SUB  R4, R4, 1
  JNZ  R4, loop
\end{lstlisting}

\begin{longtable}{L{0.18\linewidth}L{0.38\linewidth}L{0.34\linewidth}}
\toprule
程序结构 & PC 动作 & 硬件证据 \\
\midrule
顺序指令 & \code{PC = PC + size} & 取指地址递增，\code{pc_sel=sequential} \\
条件分支 & 根据标志选择顺序 PC 或目标 PC & 标志位、条件译码、目标加法器 \\
循环回边 & 分支目标小于当前 PC 附近地址 & 同一组编码被重复取指 \\
函数调用 & 保存返回地址并选择函数入口 & 链接寄存器或栈写入，PC 改为目标 \\
返回 & 从寄存器或栈恢复 PC & \code{RET} 或间接跳转选择返回地址 \\
\bottomrule
\end{longtable}

用这个视角看循环，性能问题也会具体起来。循环体中的 LOAD 可能 cache miss，ADD 依赖 LOAD 结果，分支依赖计数器更新，取指路径又要处理回跳。简单 CPU 会逐条完成；流水线 CPU 需要处理 load-use 冒险和分支控制冒险；带预测的 CPU 会猜测回跳是否成立。无论复杂度如何，循环仍然是 PC、寄存器、标志和内存状态反复经过同一批硬件路径。

\topic{调用约定把寄存器和栈的责任写成合同}
硬件只提供寄存器、LOAD/STORE、PC 改写和少数专门指令；函数调用是否能可靠嵌套，还需要调用约定。调用约定规定参数放在哪些寄存器或栈位置，返回值放在哪里，哪些寄存器由调用者保存，哪些由被调用者保存，栈指针如何对齐。它不是纯软件礼节，因为每一条规则都会变成具体寄存器写回和内存访问。

\begin{longtable}{L{0.2\linewidth}L{0.36\linewidth}L{0.34\linewidth}}
\toprule
调用约定对象 & 机器级作用 & 硬件动作 \\
\midrule
参数寄存器 & 前几个参数直接进入寄存器 & 调用前写寄存器，被调用函数读寄存器 \\
返回值寄存器 & 函数结果放在约定寄存器 & 返回前写回，调用者随后读取 \\
调用者保存寄存器 & 调用后不保证保持旧值 & 调用者必要时先 STORE 到栈 \\
被调用者保存寄存器 & 函数返回前必须恢复 & 被调用者入口 STORE，出口 LOAD \\
栈对齐 & 保证局部对象、返回地址和向量访问边界 & 调整 SP，可能插入填充字节 \\
\bottomrule
\end{longtable}

\begin{lstlisting}
caller:
  STORE R5, [SP - 4]   // caller saves if it needs R5 later
  SP = SP - 4
  R1 = arg0
  R2 = arg1
  CALL target
  SP = SP + 4
  LOAD R5, [SP - 4]
  use R0               // return value

callee:
  SP = SP - frame_size
  STORE R6, [SP + 0]   // callee-saved register
  ...
  LOAD R6, [SP + 0]
  SP = SP + frame_size
  RET
\end{lstlisting}

这个 trace 能帮助读者理解为什么栈损坏常常表现成“返回后跑飞”。返回地址、保存寄存器和局部数组可能同在一段连续内存中。若局部数组越界 STORE，覆盖了保存寄存器，错误会在函数返回后才显现；若覆盖了返回地址，RET 会把被污染的字节装入 PC。CPU 执行 RET 时并不知道这个地址是否来自真实 CALL，它只执行“从约定位置取一个地址并写入 PC”的硬件动作。

\topic{异常和中断是受控的非顺序 PC 更新}
异常、中断和系统调用都可以视为特殊控制流。普通分支由指令显式请求；异常由当前指令执行过程中发现的条件触发，例如非法 opcode、除零、未对齐访问、页错误或总线错误；中断来自外部设备或定时器。三者最终都要让 CPU 保存足够的旧状态，并把 PC 改到一个规定入口。

\begin{longtable}{L{0.18\linewidth}L{0.38\linewidth}L{0.34\linewidth}}
\toprule
事件 & 触发来源 & 入口动作 \\
\midrule
同步异常 & 当前指令导致，如非法指令、页错误、除零 & 保存故障 PC 和原因，跳异常向量 \\
外部中断 & 设备、定时器或中断控制器请求 & 在可接收边界保存 PC/Flags，跳中断入口 \\
系统调用 & 指令显式请求进入特权服务 & 切换权限和入口，保存返回位置 \\
复位 & 电源或复位逻辑强制初始化 & PC 装入复位向量，控制状态清零或设默认值 \\
\bottomrule
\end{longtable}

异常路径最重要的是提交边界。若一条 LOAD 因页错误失败，目标寄存器不应保存随机数据；若一条 STORE 访问设备前发生权限错误，设备不应看到写命令；若一条除法发生除零异常，后续指令不应像已经成功执行那样提交。简单 CPU 可以在每条指令结束时才接收中断；复杂流水线 CPU 则需要记录哪些指令已经可以提交，哪些仍可能被冲刷。

\begin{lstlisting}
faulting LOAD:
  IF:  fetch instruction
  ID:  decode LOAD
  EX:  form address
  MEM: address translation fails
  WB:  suppressed
  PC:  exception_vector
\end{lstlisting}

这个例子说明，异常不是“多跳一次分支”这么简单。它还要阻止错误写回，把失败原因带给入口，并保证重试或终止时系统能判断发生了什么。第六章讨论 80386 页错误时，会看到硬件如何保存故障线性地址和错误码；第四章这里先建立统一模型：异常是硬件检测到无法继续普通提交后，选择另一条受控 PC 更新路径。

\topic{精确异常要求程序员可见状态像顺序执行}
流水线和乱序执行会让多条指令同时处在不同阶段。若较晚指令已经写回，较早指令随后发现异常，程序员可见状态就会变得难以恢复。精确异常要求异常发生时，状态看起来像按程序顺序执行到某条指令：异常之前的指令都已提交，异常指令和之后的指令都没有提交。教学 CPU 可以通过按序完成天然满足；现代 CPU 需要提交队列、重排序缓冲或类似机制。

\begin{longtable}{L{0.18\linewidth}L{0.38\linewidth}L{0.34\linewidth}}
\toprule
机制 & 解决的问题 & 代价 \\
\midrule
按序流水线 & 后一条指令不越过前一条提交 & 吞吐有限，长延迟指令会拖住后续 \\
提交队列 & 执行可乱序，提交按程序顺序 & 需要记录目的寄存器、异常原因和结果 \\
寄存器重命名 & 消除部分假依赖，允许更多并行 & 需要映射表、空闲表和恢复机制 \\
冲刷流水线 & 分支预测错或异常时丢弃错误路径 & 浪费已经做的工作，需要恢复 PC 和状态 \\
\bottomrule
\end{longtable}

虽然本书不展开乱序核心设计，但提前理解“执行”和“提交”的区别很重要。ALU 算出结果只是执行完成，结果进入程序员可见寄存器才是提交。LOAD 从 cache 取回数据只是访存完成，目标寄存器被写回才是提交。STORE 更复杂：它可能先进入 store buffer，等确认没有异常、顺序允许、地址属性允许后才对 cache、内存或设备可见。这个区分把第四章 CPU trace、第五章内存顺序和第六章平台异常连成一条线。

\topic{最小整机实验要先有可观察边界}若把本节落到一块低速实验板，不必立即追求完整 ISA。可以先定义四条动作：\code{ADD}、\code{LOAD}、\code{STORE}、\code{JZ}。指令 ROM 提供编码，4 个寄存器保存通用值，74HC283/74HC86 构成加减核心，SRAM 或寄存器阵列模拟数据存储器，若干 LED 和拨码开关映射成 MMIO。验收重点是每一步都可观察：PC LED 显示当前取指地址，IR LED 或逻辑分析仪显示当前指令，控制线可单步观察，写回和 PC 更新只在时钟沿发生。

\begin{longtable}{L{0.18\linewidth}L{0.38\linewidth}L{0.34\linewidth}}
\toprule
实验对象 & 可用器件或实现 & 验收证据 \\
\midrule
指令 ROM & EEPROM、Flash 或预置拨码/仿真 ROM & 复位后固定地址输出第一条编码 \\
PC/IR & 74HC161、74HC574 或触发器组 & PC 每步按控制更新，IR 保存本周期指令 \\
数据 RAM & SRAM、寄存器阵列或仿真存储器 & LOAD/STORE 地址、数据、写使能可观察 \\
地址译码 & 74HC138、比较器或门阵列 & RAM、LED、按键寄存器不会同时响应 \\
MMIO LED & 锁存器驱动 LED & 写指定地址后 LED 状态在边界改变 \\
MMIO 按键 & 上拉/下拉、施密特、同步器 & 读指定地址返回已同步状态，不直接用生按键 \\
\bottomrule
\end{longtable}

\topic{最小整机的单步调试顺序}搭建教学 CPU 时，最可靠的调试顺序不是一次性运行程序，而是把每个状态边界单独验收。先让复位把 PC 固定到入口地址，再单步确认 ROM 输出第一条指令；随后确认 IR 能锁存该指令，译码能产生正确控制线；再用一条 ADD 验证寄存器读、ALU 和写回；最后才加入 LOAD、STORE、分支和 MMIO。

\begin{longtable}{L{0.16\linewidth}L{0.38\linewidth}L{0.36\linewidth}}
\toprule
调试阶段 & 操作 & 通过标准 \\
\midrule
复位 & 拉低/拉高复位并给一个时钟 & PC 固定为入口，写使能均处于安全值 \\
取指 & 单步一个取指周期 & ROM 片选有效，IR 保存预期编码 \\
ADD & 预置 R1/R2，执行 \code{ADD R3,R1,R2} & R3 在时钟沿变为正确和，PC 前进 \\
LOAD & RAM 某地址预置数据，执行加载 & 地址、读使能、写回来源均正确 \\
STORE & 执行写 RAM 或写 LED 寄存器 & 只有目标片选有效，非目标不被改写 \\
分支 & 改变 Zero 条件重复执行 & PC 在顺序地址和分支目标之间正确选择 \\
\bottomrule
\end{longtable}

这种验收方式也适用于 FPGA。逻辑分析仪、仿真波形或片上 ILA 应至少抓取 PC、IR、控制字、ALU 输入输出、内存地址、读写使能、写回数据和寄存器写地址。若只看最终 LED 是否闪烁，错误可能被循环掩盖；若能看到每条指令的提交边界，错误会落到具体周期。

\topic{流水线把吞吐率问题提前暴露出来}本节及其后的 scoreboard 专题属于扩写目录中“流水线、冒险和异常”一章的提前浓缩，标记为选读：第一遍读最小 CPU 时可以先跳过，等单周期 trace 走通后再回来。一旦把取指、译码、执行、访存和写回拆成流水级，CPU 的平均吞吐率可以提高，但控制复杂度也立刻上升。相邻指令可能争用同一个存储器端口，后一条指令可能读取前一条尚未写回的寄存器，分支可能让已经取到的指令变成错误路径。这些问题统称结构冒险、数据冒险和控制冒险。

\begin{longtable}{L{0.18\linewidth}L{0.34\linewidth}L{0.38\linewidth}}
\toprule
冒险类型 & 例子 & 硬件处理 \\
\midrule
结构冒险 & 取指和 LOAD 同时争用单端口存储器 & 分离指令/数据存储，或插入停顿 \\
数据冒险 & \code{ADD R1,...} 后立刻 \code{SUB R2,R1,...} & 旁路/转发，或等待写回后再读 \\
控制冒险 & 条件分支结果未出时已取后续指令 & 预测、延迟槽、停顿或冲刷流水线 \\
异常冒险 & 较晚指令先完成，较早指令发生异常 & 提交队列或按序提交保证精确状态 \\
\bottomrule
\end{longtable}

因此，现代芯片低时延不只是“流水线更深”。流水线必须配套转发网络、停顿控制、冲刷控制和精确异常边界。否则，吞吐率提高会以错误状态提交为代价。学习最小 CPU 时提前看到这些冒险，可以防止把流水线误解成简单地把框图切成几段。

\topic{案例：从一个计数循环走到逐周期 trace}
为了把本章从概念推进到可验收工程，可以完整走读一个小程序。假设内存中从 \code{BASE} 开始有若干 32-bit 数，R1 保存当前地址，R4 保存剩余元素个数，R3 保存累加和。程序每轮加载一个元素，加入累加和，地址加 4，计数减 1，若计数不为 0 就回到循环入口。

\begin{lstlisting}
loop:
  LOAD R2, [R1 + 0]
  ADD  R3, R3, R2
  ADDI R1, R1, 4
  SUBI R4, R4, 1
  JNZ  R4, loop
done:
  STORE R3, [R0 + RESULT]
\end{lstlisting}

这段程序覆盖了本章大部分通路：取指、译码、地址形成、数据 LOAD、ALU 加法、立即数扩展、条件码或比较、条件分支、STORE。若它能在单步下跑通，最小 CPU 的核心路径基本都被验证过。

\begin{longtable}{L{0.14\linewidth}L{0.28\linewidth}L{0.32\linewidth}L{0.16\linewidth}}
\toprule
指令 & 主要输入 & 组合路径 & 提交对象 \\
\midrule
\code{LOAD} & R1、偏移 0 & ALU 地址加法，数据存储器读 & R2，PC \\
\code{ADD} & R3、R2 & ALU 加法，写回 mux & R3，Flags，PC \\
\code{ADDI} & R1、立即数 4 & 立即数扩展，ALU 加法 & R1，Flags，PC \\
\code{SUBI} & R4、立即数 1 & B 取反加一或减法路径 & R4，Zero，PC \\
\code{JNZ} & R4 或 Zero 标志 & 条件判断，PC mux & PC \\
\code{STORE} & R3、RESULT 地址 & 地址形成，数据存储器写 & 内存或 MMIO，PC \\
\bottomrule
\end{longtable}

若采用五阶段流水线，前两条指令立即暴露 load-use 数据冒险。 \code{LOAD R2} 的数据通常到 MEM 阶段末尾才回来，而下一条 \code{ADD R3,R3,R2} 在 EX 阶段就需要 R2。没有转发或停顿时，ADD 会读到 R2 的旧值。

\begin{longtable}{L{0.12\linewidth}L{0.18\linewidth}L{0.18\linewidth}L{0.18\linewidth}L{0.18\linewidth}}
\toprule
周期 & LOAD & ADD & ADDI & SUBI \\
\midrule
1 & IF &  &  &  \\
2 & ID & IF &  &  \\
3 & EX & ID & IF &  \\
4 & MEM & stall 或 EX 等待 & ID & IF \\
5 & WB & EX 使用转发值 & ID/EX & ID \\
\bottomrule
\end{longtable}

这个表不是为了背流水线图，而是训练读者问三个问题：数据什么时候产生，下一条什么时候需要，是否有合法路径把新值送过去。若没有旁路网络，就必须插入停顿；若有从 MEM/WB 到 EX 的转发，控制器必须检测依赖并选择转发输入。否则程序在低速单周期模型中正确，在流水线模型中错误。

\topic{案例：条件分支和预测错误怎样恢复}
同一个循环的 \code{JNZ} 还能展示控制冒险。取指单元通常希望每周期都取下一条指令，但分支是否成立要等比较结果出来。若等到结果再取指，流水线会空等；若提前猜测，就可能取错路径。教学 CPU 可以先选择“遇到分支就停顿”，再讨论预测。

\begin{longtable}{L{0.18\linewidth}L{0.38\linewidth}L{0.34\linewidth}}
\toprule
策略 & 硬件动作 & 代价 \\
\midrule
停顿直到分支结果 & 分支进入 EX 前不继续提交后续路径 & 控制简单，循环每次损失周期 \\
总是假设不跳 & 继续取顺序 PC，若实际跳转则冲刷 & 适合少跳分支，循环回边常错 \\
总是假设跳转 & 提前取目标 PC，若实际不跳则冲刷 & 循环体内常较好，退出时错一次 \\
动态预测 & 根据历史选择方向和目标 & 需要预测表、目标缓存和恢复机制 \\
\bottomrule
\end{longtable}

预测错误恢复必须同时处理两类状态。第一，已经取到但不该执行的指令要被冲刷，不能写寄存器、内存或设备。第二，PC 要恢复到正确目标，取指从那里重新开始。若错误路径上有 STORE 或 MMIO 写已经对外可见，就无法简单恢复；因此流水线通常要把真正对外可见的提交推迟到确认顺序正确之后。

\begin{lstlisting}
branch predicted not taken:
  fetch sequential instruction
  branch resolves taken
  squash wrong-path instruction
  PC = branch_target
  restart fetch
\end{lstlisting}

\topic{案例：函数调用的完整栈帧 trace}
再用一个更接近系统程序的例子。函数 \code{sum2(a,b)} 返回两个整数之和。调用者把参数放进 R1/R2，CALL 保存返回地址并跳转；被调用者可能建立小栈帧，保存需要保留的寄存器，计算结果放进 R0，RET 恢复 PC。

\begin{lstlisting}
caller:
  LDI  R1, 7
  LDI  R2, 5
  CALL sum2
  STORE R0, [R0 + RESULT]

sum2:
  SP = SP - 4
  STORE R5, [SP + 0]
  ADD R0, R1, R2
  LOAD R5, [SP + 0]
  SP = SP + 4
  RET
\end{lstlisting}

\begin{longtable}{L{0.18\linewidth}L{0.38\linewidth}L{0.34\linewidth}}
\toprule
调用阶段 & 程序员可见状态变化 & 硬件边界 \\
\midrule
准备参数 & R1=7，R2=5 & LDI 写回寄存器，PC 顺序前进 \\
CALL & 返回地址保存，PC=sum2 入口 & 写链接寄存器或 STORE 到栈，再写 PC \\
建立栈帧 & SP 减小，保存 R5 & ALU 改 SP，STORE 到栈地址 \\
计算返回值 & R0=R1+R2 & ALU 加法，R0 写回 \\
销毁栈帧 & 恢复 R5，SP 加回 & LOAD、ALU 加法、寄存器写回 \\
RET & PC=返回地址 & 从链接寄存器或栈读地址，写 PC \\
\bottomrule
\end{longtable}

这个例子能把“调用约定”和“硬件状态”绑在一起。若 SP 初始化错误，第一条保存寄存器的 STORE 就可能写坏数据区；若 CALL 保存返回地址的位置错误，RET 会跳到不可预测地址；若被调用者忘记恢复 R5，调用者看到的 R5 就被破坏。所谓 ABI 兼容，最终表现为这些寄存器、栈位置和 PC 更新都按同一合同执行。

\topic{案例：异常插入到调用链中}
若 \code{sum2} 内部执行 LOAD 时发生页错误或总线错误，异常入口必须知道故障发生在哪条指令、如何保存返回上下文、是否能恢复。简单裸机可能直接停机或点亮错误 LED；带操作系统的机器会把故障地址、错误码和当前 PC 保存到异常栈，再进入处理程序。

\begin{longtable}{L{0.18\linewidth}L{0.42\linewidth}L{0.3\linewidth}}
\toprule
异常保存对象 & 为什么需要 & 恢复或诊断用途 \\
\midrule
故障 PC & 知道哪条指令失败 & 修复后重试，或报告故障位置 \\
状态/标志 & 保留中断允许、条件码、特权状态 & 返回时恢复执行环境 \\
通用寄存器 & 处理程序可能使用寄存器 & 避免破坏被打断程序 \\
错误原因 & 页不存在、权限错误、总线超时等 & 选择缺页、杀进程、复位设备或停机 \\
栈指针/特权栈 & 异常处理本身需要可靠栈 & 防止在坏用户栈上继续保存状态 \\
\bottomrule
\end{longtable}

因此，异常处理不是在程序旁边加一条旁路，而是 CPU、栈、向量表、权限和内存映射共同定义的第二套控制流。第四章到这里可以停在一个结论：最小 CPU 只要能顺序执行还不够，真正的机器还必须能在错误、等待、分支和外部事件中保持可恢复的状态边界。

\topic{本章到此只验收 CPU 内部 trace}最小 CPU 章节到这里应停在内部状态边界：PC、IR、寄存器堆、ALU、写回 mux、控制字、等待状态和流水线冒险。只要读者能把每条指令写成“读旧状态、经过组合路径、在时钟沿提交新状态”的 trace，本章目标就已经完成。

下一章再把这个 CPU 接到主存、总线、MMIO、DMA 和中断，讨论地址译码、等待、错误、cache 可见性和设备副作用。第六章再用 6502、Z80、8051、8086、80386、80486/Pentium 和现代 SoC 观察真实芯片如何把这些边界逐步扩展到整机平台。这样分章后，CPU 内部 trace、整机事务、经典芯片三类问题不会混在同一节里。

\topic{专题：译码器怎样从指令字段产生控制线}
指令不是被 CPU “理解”后才执行，而是字段被译码成一组控制选择。操作码决定大类，寄存器字段连接寄存器堆读写地址，立即数字段进入符号扩展或零扩展路径，功能字段进一步选择 ALU 操作。最小 CPU 的译码器可以是组合逻辑，也可以是 ROM/PLA 或微码入口，但它的产物都应落到可检查的控制线上。

\begin{longtable}{L{0.18\linewidth}L{0.34\linewidth}L{0.38\linewidth}}
\toprule
指令字段 & 进入的硬件路径 & 典型错误 \\
\midrule
\code{opcode} & 主控制器，决定 \code{mem_read/reg_write/pc_sel} 等 & 未定义 \code{opcode} 被当成合法指令执行 \\
\code{rs/rt} & 寄存器堆读地址 & 字段位置接反导致读错源寄存器 \\
\code{rd} & 寄存器堆写地址或写回 mux 的目标选择 & STORE 或 BRANCH 错误写寄存器 \\
\code{imm} & 符号扩展、零扩展、左移、地址加法器 & 立即数扩展规则与 ISA 不一致 \\
\code{funct} & ALU 控制的二级译码 & ADD/SUB/SLT 等同 \code{opcode} 下混淆 \\
\bottomrule
\end{longtable}

译码必须明确非法指令策略。教学机器可以停机并点亮错误码；带异常的机器可以保存 PC 并进入非法指令入口；极简硬件也可以规定未定义 opcode 为 NOP，但这种做法会掩盖程序损坏。无论选择哪种策略，都不能让非法组合悄悄打开 \code{mem_write} 或跳到不可预测 PC。

\topic{专题：scoreboard 和冒险检测的最小逻辑（选读）}本专题与前面的流水线冒险同属“流水线、冒险和异常”一章的提前浓缩，选读。流水线把多条指令重叠执行后，控制器必须知道“某个结果现在是否已经可用”。最简单的办法不是先设计复杂乱序执行，而是在流水级寄存器里保留目标寄存器号、是否写回、结果来自 ALU 还是存储器等信息，再由冒险检测逻辑比较当前指令的源寄存器。

\begin{longtable}{L{0.18\linewidth}L{0.44\linewidth}L{0.28\linewidth}}
\toprule
冒险 & 检测证据 & 处理策略 \\
\midrule
ALU-ALU RAW & 前一条会写 \code{rd}，后一条读取同一寄存器 & 转发 EX/MEM 或 MEM/WB 结果 \\
load-use & LOAD 的目标寄存器被下一条立即读取 & 停顿一拍，等待数据从存储阶段回来 \\
STORE 数据 & STORE 的写数据来自尚未写回的指令 & 转发到存储数据输入，或停顿 \\
分支控制 & 分支结果尚未决定但后续指令已取入 & 冲刷流水线中错误路径指令 \\
结构冲突 & 取指和访存争用同一存储端口 & 停顿，或分离指令/数据存储器 \\
\bottomrule
\end{longtable}

最小冒险检测可写成硬件式条件，而不是口头描述。例如：
\[
\text{stall} = \text{IDEX.memRead}\land((\text{IDEX.rd}=\text{IFID.rs})\lor(\text{IDEX.rd}=\text{IFID.rt}))
\]
这表示当前处于执行级的 LOAD 还没有数据，而取指/译码级指令马上要读这个目标寄存器。停顿时 PC 和 IF/ID 不更新，ID/EX 注入 bubble，避免错误指令带着旧数据前进。

\topic{专题：同一条 LOAD 在三种 CPU 组织中的对照}
为了避免“单周期、多周期、流水线”只停留在名词层面，可以用同一条 \code{LOAD R1,[R2+8]} 对照三种组织。它们完成的语义相同：读 R2，加立即数形成地址，访问存储器，把数据写回 R1。区别在于资源是否复用、状态保存在哪里、关键路径由谁决定。

\begin{longtable}{L{0.16\linewidth}L{0.42\linewidth}L{0.32\linewidth}}
\toprule
组织方式 & LOAD 的执行形态 & 主要代价 \\
\midrule
单周期 & 取指、译码、寄存器读、地址加法、存储读、写回全部在一拍完成 & 时钟周期被最慢指令拉长 \\
多周期 & IF、ID、EX 地址形成、MEM、WB 分拍完成 & 控制状态机复杂，但硬件可复用 \\
流水线 & 不同 LOAD 位于不同阶段，阶段间用寄存器隔开 & 需要处理冒险、冲刷和阶段平衡 \\
\bottomrule
\end{longtable}

单周期设计最适合教学入门，因为一条指令的 trace 很直观；多周期设计更像早期现实 CPU 的控制方式，因为它允许慢步骤独占多个周期；流水线设计提高吞吐量，但不会让单条 LOAD 的物理工作消失。读者应能把三种设计画成状态边界，而不是只背“流水线更快”。

\topic{专题：控制器验收要同时看正路径和错误路径}
CPU 章节的验收不能只跑一段加法程序。控制器必须面对非法 opcode、未对齐访问、慢存储器、分支冲刷、中断屏蔽、异常返回和复位。否则机器在正常程序上看似通过，一旦接入真实总线或设备就会暴露状态边界不清的问题。

\begin{longtable}{L{0.2\linewidth}L{0.42\linewidth}L{0.28\linewidth}}
\toprule
场景 & 必须观察的状态 & 合格表现 \\
\midrule
非法 opcode & PC、IR、错误码或 halt 状态 & 不写内存，不错误提交寄存器 \\
慢 LOAD & 地址、\code{mem_read}、\code{ready}、MDR & 等待期间请求稳定，只提交一次 \\
分支命中 & \code{IF/ID}、\code{ID/EX}、PC 选择 & 错误路径指令被冲刷 \\
异常入口 & 故障 PC、状态寄存器、异常向量 & 可诊断、可恢复或明确停机 \\
复位 & PC、寄存器堆、控制状态机 & 无旧控制线残留，第一条取指确定 \\
\bottomrule
\end{longtable}

这一类验收会迫使实现者回答：哪个周期算提交？哪些控制线在等待时保持？错误路径是否会产生副作用？这些答案比“程序最后算出 42”更接近 CPU 工程。



\section*{章末练习与验收任务}

本章练习不要求先追求完整 ISA，而要求把每条指令写成可检查的硬件 trace。答案必须同时说明控制线、数据路径和提交边界。

\begin{longtable}{L{0.2\linewidth}L{0.48\linewidth}L{0.22\linewidth}}
\toprule
任务 & 要完成的产物 & 验收标准 \\
\midrule
PC 取指 & 写出 PC、指令存储器、IR 和 \code{pc_write/ir_write} 的周期关系 & 能说明哪一拍保存指令 \\
ADD trace & 为 \code{R3=R1+R2} 列出读地址、ALU 操作、写回选择和写使能 & 写回只在目标时钟沿发生 \\
LOAD trace & 把 LOAD 拆成地址形成、存储读、MDR 保存、寄存器写回 & 能处理存储器未 ready \\
STORE trace & 说明地址、写数据、字节使能、写控制和无寄存器写回 & 能区分 RAM 写和 MMIO 副作用 \\
JZ trace & 说明 Zero 标志如何影响下一 PC & 分支目标和顺序 PC 选择明确 \\
条件码比较 & 分别解释有符号小于和无符号小于使用哪些 ALU 证据 & 不把 Carry、Sign 和 Overflow 混用 \\
CALL/RET trace & 把调用和返回拆成保存返回地址、更新 SP、改写 PC 和恢复 PC & 能指出栈访问本质上是 LOAD/STORE \\
控制字 & 设计一个包含 \code{alu_op/wb_sel/reg_write/mem_read/mem_write/pc_sel} 的控制字 & 每一位能追到实际电路 \\
微程序 & 为 LOAD 写出 5 个微步骤和等待状态处理 & 慢存储器只拖住访存阶段 \\
流水线冒险 & 说明结构、数据、控制三类冒险各破坏什么边界 & 能提出停顿、转发或冲刷策略 \\
\bottomrule
\end{longtable}

\section*{参考解答与验收要点}

\begin{longtable}{L{0.18\linewidth}L{0.5\linewidth}L{0.22\linewidth}}
\toprule
任务 & 参考解答 & 必须保留的验收点 \\
\midrule
PC 取指 & PC 输出取指地址，指令存储器返回编码，\code{ir_write} 在取指边界保存编码；下一 PC 可由 \code{pc+size}、分支目标或异常入口选择。 & PC 和 IR 都是状态元件。 \\
ADD trace & 读端口选择 R1/R2，ALU 选择 ADD，写回 mux 选择 ALUOut，写地址为 R3，\code{reg_write=1} 只在写回边界有效。 & 算对不等于写对。 \\
LOAD trace & EX 阶段形成地址，MEM 阶段发起读，目标未 ready 时保持地址和控制，响应后保存 MDR，WB 阶段写回目标寄存器。 & 等待不能污染其他状态。 \\
STORE trace & STORE 使用 ALU 形成地址，源寄存器提供写数据，打开字节使能和写控制，不打开寄存器写回。若地址命中 MMIO，副作用由设备规格决定。 & STORE 是有副作用的事务。 \\
JZ trace & 比较或 ALU 结果产生 Zero，控制器根据 Zero 选择顺序 PC 或分支目标；PC 只在提交边界更新。 & 条件输入必须在采样前稳定。 \\
控制字 & 控制字字段应覆盖 PC、IR、寄存器堆、ALU、存储器和下一状态逻辑；任何字段没有连接对象都不是硬件控制。 & 控制字是电线集合，不是注释。 \\
微程序 & LOAD 可写为取指、读寄存器、地址形成、访存等待、写回；访存未完成时微地址保持在等待状态，错误时跳异常入口。 & 微步骤提供可观察边界。 \\
流水线冒险 & 结构冒险来自资源争用，数据冒险来自结果未写回，控制冒险来自下一 PC 尚未确定。停顿保护边界，转发减少等待，冲刷丢弃错误路径。 & 提高吞吐不能破坏提交顺序。 \\
\bottomrule
\end{longtable}

\begin{classicnote}
最小 CPU 的经典读法是逐周期 trace。不要只画“PC 连到内存、ALU 连到寄存器”的框图；要写清楚每条指令在哪个周期读旧状态、在哪条组合路径上传播、在哪个时钟沿提交新状态。后续主存、总线、I/O 和经典芯片只是在这个 trace 外面继续增加事务边界。
\end{classicnote}
